\subsection{Spin-Lifted Deficit Factor}
The non-orientable identification acts as a $\mathbb{Z}_2$ operation on the base,
but in the spin bundle the same element lifts to order $4$ because a $2\pi$ rotation
acts as $-1$ on spinors and only a $4\pi$ rotation is trivial. This produces a
universal deficit factor
\begin{equation}
    1-\frac{1}{4}=\frac{3}{4},
\end{equation}
per minimal sector.

\subsection{Sector Count from Cut Geometry}
Cutting along $\Gamma$ yields boundary components $(C_1,C_2,\Gamma_+,\Gamma_-)$.
The four local corner patches
$(C_1,\Gamma_+)$, $(C_1,\Gamma_-)$, $(C_2,\Gamma_+)$, $(C_2,\Gamma_-)$
constitute the fundamental domain. Summing the spin-lifted deficit factor over these sectors yields
\begin{equation}
    \sum_{p=1}^{F}\left(1-\frac{1}{4}\right)=4\cdot\frac{3}{4}=3.
\end{equation}
In TFPT normalization this yields the canonical form $\deltatop = 3/(256\pi^4)$.

\subsection{\texorpdfstring{Conversion to the $\cthree$ normalization}{Conversion to the c3 normalization}}

Using $\cthree = 1/(8\pi)$:
\begin{equation}
    48\cthree^4 = 48 \times \frac{1}{(8\pi)^4} 
    = \frac{48}{4096\pi^4}
    = \frac{3}{256\pi^4}
\end{equation}
so $\deltatop = 3/(256\pi^4)$ and $\deltatop = 48\cthree^4$ are exactly equivalent.

\subsection{Status note}
The spin-lift argument fixes the $1-1/4$ deficit factor robustly, and the sector count
is fixed by seam symmetry. The seam normalization is provided by the $\eta$-gluing
derivation in Appendix~\ref{app:seam}.

%=============================================================================
\section{Flavor Monodromy and Cross-Ratio}
\label{app:flavor}
%=============================================================================

This appendix derives the M\"obius structure of flavor relations from boundary monodromy.

\subsection{Projective Invariance on Boundaries}

On a conformally flat boundary, the metric is locally:
\begin{equation}
    ds^2 = e^{2\omega(x)} \delta_{ij} dx^i dx^j
\end{equation}
Holomorphic automorphisms (Wilson lines of flat connections) act as M\"obius 
transformations:
\begin{equation}
    z \mapsto \frac{az + b}{cz + d}, \quad ad - bc = 1
\end{equation}

\begin{remark}
\textbf{Context:} $\mathrm{SL}(2,\mathbb{C})$ double-covers the proper orthochronous Lorentz group $\mathrm{SO}^+(3,1)$, 
and its projectivization $\mathrm{PSL}(2,\mathbb{C})$ acts by M\"obius transformations on the Riemann sphere. 
This links the M\"obius structure to standard Lorentz symmetry rather than an ad hoc choice.
\end{remark}

\subsection{Cross-Ratio as Unique Invariant}

Given four marked points $\{z_1, z_2, z_3, z_4\}$ on the boundary, the unique 
projective invariant is the cross-ratio:
\begin{equation}
    [z_1, z_2; z_3, z_4] = \frac{(z_1 - z_3)(z_2 - z_4)}{(z_1 - z_4)(z_2 - z_3)}
\end{equation}

\textbf{Consequence:} Any physical quantity depending only on boundary data must be 
expressible as a function of cross-ratios, hence as a M\"obius map.

\subsection{Cusp Selection}

The cusps $y \in \{1, 1/3, 2/3\}$ arise from:
\begin{enumerate}[leftmargin=1.5em, topsep=2pt, itemsep=1pt]
    \item \textbf{Rationality:} Quantized holonomies $\Rightarrow$ rational cusps.
    \item \textbf{GUT compatibility:} SU(5) hypercharge assignments give 
        $Y = 1/3$ (quarks), $Y = 1$ (leptons), etc.
    \item \textbf{Uniqueness:} These are the only rational values that simultaneously 
        separate the three SM sectors (up, down, lepton) and maintain GUT normalization.
\end{enumerate}

\subsection{\texorpdfstring{Derivation of $\delta_\star$}{Derivation of delta_star}}

The flavor phase is a first-order closure:
\begin{equation}
    \delta_\star = \delta_0 + \delta_1 \cdot \phiz + \mathcal{O}(\phiz^2)
\end{equation}

\textbf{Leading term:} $\delta_0 = 3/5$ from the abelian trace (same as $\bone$ numerator 
divided by 41/10's denominator structure).

\textbf{Correction:} $\delta_1 = 1/6$ from the M\"obius length scale. The factor $1/6$ 
matches the boundary normalization ($6\pi$ total curvature).

\textbf{Result:}
\begin{equation}
    \delta_\star = \frac{3}{5} + \frac{\phiz}{6} = 0.6 + 0.00886 = 0.6089
\end{equation}

%=============================================================================
\section{Reproducibility}
\label{app:repro}
%=============================================================================

\textbf{Replication package:} All numerical results can be reproduced using the 
TFPT verification suite (v5).

\textbf{Contents:}
\begin{itemize}[leftmargin=1.5em, topsep=2pt, itemsep=1pt]
    \item Python modules for each calculation domain
    \item JSON files with all constants and derived quantities
    \item Automated test suite with pass/fail status
    \item Plot generation scripts
    \item Verification suite modules include: \texttt{effective\_action\_r2}, \texttt{bounce\_perturbations} (see \codepath{tfpt-suite/} and the generated \codepath{tfpt-suite/tfpt-test-results.pdf})
\end{itemize}

\textbf{Key output files:}
\begin{itemize}[leftmargin=1.5em, topsep=2pt, itemsep=1pt]
    \item \texttt{tfpt\_complete\_analysis.json}: Master results file
    \item \texttt{e8\_chain\_analysis.json}: Complete E$_8$ chain data
    \item \texttt{gauge\_couplings.csv}: Two-loop RGE output (via PyR@TE)
\end{itemize}

\textbf{Execution:} \texttt{python -m tfpt\_verification5.run\_all}

\textbf{Verification:} The core modules must report ``success'' status; the
$R^2$ effective-action and bounce-perturbation modules are implemented in the shipped suite and included in the generated test report.

%=============================================================================
\section{Prediction Ledger and Falsification}
\label{app:prediction-ledger}
%=============================================================================

\textbf{Purpose:} Provide a single, reviewer-friendly list of predictions and candidates that is tied to the same JSON sources used by the verification suite, to prevent number drift between Abstract, theorems, tables, and appendices.

\begingroup
\footnotesize
\renewcommand{\arraystretch}{1.15}

% AUTO-GENERATED by tfpt_suite/export_prediction_ledger.py — inlined for standalone-paper build.
\begin{longtable}{@{}p{0.15\textwidth}p{0.27\textwidth}rrrp{0.10\textwidth}p{0.18\textwidth}@{}}
\caption{Prediction ledger (auto-generated from TFPT suite outputs).}\label{tab:prediction-ledger}\\
\toprule
Observable & Formula / dependencies & Pred. & Ref. & $z$ & Status & Suite evidence\\
\midrule
\endfirsthead
\toprule
Observable & Formula / dependencies & Pred. & Ref. & $z$ & Status & Suite evidence\\
\midrule
\endhead
\midrule
\multicolumn{7}{r}{\small Continued on next page.}\\
\midrule
\endfoot
\bottomrule
\endlastfoot
\multicolumn{7}{@{}l}{\textbf{Candidates / open interfaces}}\\
\(\Omega_b\) & \(\Omega_b=(4\pi-1)\beta_{\rm rad}\) under explicit sector-counting assumptions [Candidate] & \num{0.048940662665450112200545653760944651480054952073107413870556556782266783965399009} & \num{0.049301692328524438476424754994611856173233055557122787616860658651215012327847394} $\pm$ \num{0.00085681142747140931509835404112378635706093499964689237552400217734754363189007356} & \num{-0.421364} & C (conditional) & \texttt{omega\_b\_conjecture\_scan}\\ \codepath{out/omega_b_conjecture_scan/results.json}\\
\(\rho_\Lambda\) (normalization audit) & \(\phi_\*\!=n\,e^{-\alpha^{-1}(0)/2}\), \(\rho_\Lambda=(\bar M_{\rm Pl}\phi_\*)^4\); best \(n=1/2\) with log10 mismatch [Candidate] & \num{2.060247122268795528691070560351946768145020625187219242978373038664026793360454e-47} & \num{2.514176465474032e-47} & --- & C (PASS audit) & \texttt{dark\_energy\_exponential\_audit}\\ \codepath{theoryv3/out/dark_energy_exponential_audit/results.json}\\
\(\rho_\Lambda\) (torsion condensate closure) & Gap-equation completion with discrete \(n\) (spectral flow); compare to cosmology target [Candidate] & \num{4.1204942445375910573821411207038935362900412503744384859567460773280535867209079e-47} & \num{2.514176465474032e-47} & \num{2.1403530198484351956856249171986807549266690574189652264117407874776932146992054} & C (FAIL gate) & \texttt{torsion\_condensate}\\ \codepath{out_physics/torsion_condensate/results.json}\\
\multicolumn{7}{@{}l}{\textbf{Comparison layer}}\\
\(\overline{\alpha}^{(5)}(M_Z)^{-1}\) & QED running + matching (declared policy) from \(\alpha(0)\) [Policy] & \num{127.94051870737564248749202130667255420059119636590364698278411016186034311550876} & \num{127.93} $\pm$ \num{0.008} & \num{1.314838421955310936502663334069275073899545737955872848013770232542889438594741} & comparison & \texttt{global\_consistency\_test}\\ \codepath{out/global_consistency_test/results.json}\\
\multicolumn{7}{@{}l}{\textbf{Conditional dynamics}}\\
\(A_s\) & Starobinsky \(R^2\) completion: \(A_s\simeq \frac{N^2}{24\pi^2}(M/\bar M_{\rm Pl})^2\) [Conditional] & \num{0.000000002090191364329459982659058205552757672085364626939873929875219205466615549540939} & \num{0.0000000020989031673} $\pm$ \num{0.00000000002938464434} & \num{-0.29647467805764898160213003439885527348994189187882714815192466383904684207796011} & conditional & \texttt{global\_consistency\_test}\\ \codepath{out/global_consistency_test/results.json}\\
\(\Omega_a h^2\) & Axion relic accounting (scenario-dependent) [Conditional] & \num{0.12275084115028702} & --- & --- & conditional & \texttt{axion\_dm\_audit}\\ \codepath{theoryv3/out/axion_dm_audit/results.json}\\
\(\nu\) [GHz] & \(\nu_{\rm GHz}\approx 0.24179893\, m_a[\mu\mathrm{eV}]\) [Conditional] & \num{15.76402318010866} & --- & --- & conditional & \texttt{axion\_dm\_audit}\\ \codepath{theoryv3/out/axion_dm_audit/results.json}\\
\(f_a\) [GeV] & PQ cascade block (axion claim) [Conditional] & \num{88639886850.34554} & --- & --- & conditional & \texttt{axion\_dm\_audit}\\ \codepath{theoryv3/out/axion_dm_audit/results.json}\\
\(m_a\) [\(\mu\)eV] & PQ cascade block (axion claim) [Conditional] & \num{65.1947598780055} & --- & --- & conditional & \texttt{axion\_dm\_audit}\\ \codepath{theoryv3/out/axion_dm_audit/results.json}\\
\(n_s\) & Starobinsky \(R^2\) completion: \(n_s=1-\frac{2}{N}\) (benchmark \(N=56\)) [Conditional] & \num{0.96428571428571428571428571428571428571428571428571428571428571428571428571428571} & \num{0.9649} $\pm$ \num{0.0042} & \num{-0.14625850340136054421768707482993197278911564625850340136054421768707482993197288} & conditional & \texttt{global\_consistency\_test}\\ \codepath{out/global_consistency_test/results.json}\\
\(r\) & Starobinsky \(R^2\) completion: \(r=\frac{12}{N^2}\) (benchmark \(N=56\)); compared to an upper bound [Conditional] & \num{0.003826530612244897959183673469387755102040816326530612244897959183673469387755102} & < 0.036 & --- & conditional & \texttt{global\_consistency\_test}\\ \codepath{out/global_consistency_test/results.json}\\
\multicolumn{7}{@{}l}{\textbf{Core (parameter-free)}}\\
\(\alpha^{-1}(0)\) & \(\mathrm{CFE} + \varphi_0(\alpha)\) with \(\delta_2=(g/4)\delta_{\mathrm{top}}^2\) and \(g=5\) [Kernel+Grammar] & \num{137.03599921615843479026171751711397398605568500258186146747743646328176776397918} & \num{137.035999177} $\pm$ \num{0.000000021} & \num{1.8646873709648436912911416183836040477419934032132112601562746554275798808614493} & P & \texttt{defect\_partition\_g5\_audit}\\ \codepath{theoryv3/out/defect_partition_g5_audit/results.json}\\
\(\beta_{\mathrm{deg}}\) & \(\beta_{\mathrm{deg}}=\frac{180}{\pi}\frac{\varphi_0}{4\pi}\) [Kernel] & \num{0.24243503090092952849243152047951445246859512239848881592836053961070607871616298} & \num{0.35} $\pm$ \num{0.14} & \num{-0.7683212078505033679112034251463253395100348400107941719402818599235280091702644} & P & \texttt{global\_consistency\_test}\\ \codepath{out/global_consistency_test/results.json}\\
\(\lambda\equiv |V_{us}|\) & \(\lambda=\sqrt{\varphi_0}\left(1-\frac{1}{2}\varphi_0\right)\) [Kernel] & \num{0.22445997051897373245511597747748194821840297007791297390969270711773328122860166} & \num{0.22501} $\pm$ \num{0.00068} & \num{-0.80886688386215815424120959193831144352504400306915601515778365039223348735049879} & P & \texttt{global\_consistency\_test}\\ \codepath{out/global_consistency_test/results.json}\\
\(\sin^2\theta_{12}\) & TM1 sum rule from \(\sin^2\theta_{13}\) [Kernel+Grammar] & \num{0.31792771} & \num{0.307} $\pm$ \num{0.012} & \num{0.910642} & P & \texttt{pmns\_z3\_breaking}\\ \codepath{out/pmns_z3_breaking/results.json}\\
\(\sin^2\theta_{13}\) & Identity + TM1 [Kernel+Grammar] & \num{0.02310844} & \num{0.02224} $\pm$ \num{0.00057} & \num{1.523570} & P & \texttt{pmns\_z3\_breaking}\\ \codepath{out/pmns_z3_breaking/results.json}\\
\(g_{a\gamma\gamma}\) & \(g_{a\gamma\gamma}=-4c_3=-\frac{1}{2\pi}\) [Kernel] & \num{-0.15915494309189535} & --- & --- & P & \texttt{(algebraic)}\\ \codepath{paper kernel (no suite output file)}\\
\end{longtable}
\endgroup

\textbf{Falsification guide (high-level):}
\begin{itemize}[leftmargin=1.5em, topsep=2pt, itemsep=1pt]
    \item \textbf{Core (parameter-free):} Any robust shift of CODATA $\alpha^{-1}(0)$ or a decisive inconsistency in the \texttt{defect\_partition\_g5\_audit} closure would directly falsify the TFPT kernel.
    \item \textbf{Comparison layer:} A stable discrepancy in $\overline{\alpha}^{(5)}(M_Z)^{-1}$ after adopting the declared running/matching policy falsifies the paper-to-observable bridge (not the kernel itself).
    \item \textbf{Conditional dynamics:} Failures here indicate that the adopted completion/policy layer (e.g.\ $R^2$ benchmark choices, relic accounting scenario) is not closed or not correct, and must be treated as conditional rather than unconditional prediction.
    \item \textbf{Candidates / open interfaces:} These rows are explicitly not ToE-closed; a FAIL gate is reported as such and should not be framed as a resolved prediction.
\end{itemize}

%=============================================================================
\section{Identity Catalogue and Status Labels}
\label{app:identities}
%=============================================================================

This appendix labels frequently used equalities by logical status to prevent conflating definitions, algebraic rewrites, and genuine predictions.

\textbf{Legend:}
\begin{itemize}[leftmargin=1.5em, topsep=2pt, itemsep=1pt]
    \item \textbf{D}: Definition (input)
    \item \textbf{K}: Consequence of definitions (algebraic)
    \item \textbf{L}: Lemma from axioms/postulates
    \item \textbf{P}: Prediction (output; not used as input)
\end{itemize}

\begin{center}
\renewcommand{\arraystretch}{1.2}
\begin{tabular}{@{}lp{12cm}@{}}
\toprule
\textbf{Tag} & \textbf{Statement}\\
\midrule
\textbf{D} & $\beta_{\text{rad}} := \phiz/(4\pi)$.\\
\textbf{K} & $\phiz = 4\pi\beta_{\text{rad}}$.\\
\textbf{K} & $\beta_{\text{deg}} := (180/\pi)\beta_{\text{rad}}$.\\
\textbf{K} & $\gagg = -4\cthree = -1/(2\pi)$.\\
\textbf{K} & $\frac{d\beta}{d\eta} = -\frac{1}{2}\gagg\,\frac{d a_{\text{top}}}{d\eta} = 2\cthree\,\frac{d a_{\text{top}}}{d\eta}$ (from $\gagg=-4\cthree$; Theorem~\ref{thm:birefringence}).\\
\textbf{L} & $\gamma(0)=5/6$ (Definition~\ref{def:E8cascade}; discrete closure choice).\\
\textbf{L} & $\phitree = 1/(6\pi)$ (Axiom~\ref{ax:GB}; from $\eta$-gluing seam term).\\
\textbf{L} & $\deltatop = \frac{1}{256\pi^4}\sum_{p=1}^{4}\left(1-\frac{1}{4}\right)=\frac{3}{256\pi^4}$ (Axiom~\ref{ax:TC}).\\
\textbf{P} & CFE closure (two-defect, derived partition): $\alpha^{-1}_{\text{TFPT}}(0)=\num[round-mode=places,round-precision=12]{\TFPTalphaInvGOverFourNum}$ (CODATA $z=\num[round-mode=places,round-precision=3]{\TFPTalphaZGOverFourNum}$; Theorem~\ref{thm:selfconsistent}).\\
\textbf{L} & One-defect truncation (diagnostic baseline): $\alpha^{-1}(0)=\num[round-mode=places,round-precision=10]{\TFPTalphaInvSelfNum}$ (Lemma~\ref{lem:oneDefectTruncation}).\\
\textbf{P} & Cabibbo $\lambda = \sqrt{\phiz}\left(1-\frac{1}{2}\phiz\right)=0.22445997\ldots$ (Corollary~\ref{cor:Cabibbo}).\\
\textbf{P} & Birefringence $\beta_{\text{deg}} \approx 0.2424^\circ$ from $\Delta a_{\text{top}} = n\,\phiz$ (minimal $n=1$; Theorem~\ref{thm:birefringence}; optional realization in Sec.~\ref{sec:deltaa}).\\
\textbf{P} & $R^2$ inflation: $n_s=1-2/N$, $r=12/N^2$, $A_s \simeq \frac{N^2}{24\pi^2}(M/\Mpl)^2$ with $M/\Mpl=\sqrt{8\pi}\cthree^4$ (Theorem~\ref{thm:inflation}; assumptions K1--K3).\\
\textbf{P} (spec.) & $\Omega_b = (4\pi-1)\beta_{\text{rad}}$ (Table~\ref{tab:beta}).\\
\bottomrule
\end{tabular}
\end{center}

%=============================================================================
\section{\texorpdfstring{One-Loop Effective Action in Riemann--Cartan and the $R^2$ Coefficient}{One-Loop Effective Action in Riemann--Cartan and the R2 Coefficient}}
\label{app:r2}
%=============================================================================

\begin{tcolorbox}[colback=gray!5, colframe=gray!60, 
    title={\textbf{Assumptions K1--K4}}, fonttitle=\bfseries\small]
\begin{itemize}[leftmargin=1.5em, topsep=2pt, itemsep=1pt]
    \item \textbf{K1:} Field content and torsionful Lagrangian fixed by the minimal UFE sector.
    \item \textbf{K2:} Quadratic fluctuation operator is of Laplace type after background-field gauge fixing.
    \item \textbf{K3:} Heat-kernel coefficient $a_2$ captures the curvature-squared terms in $\Gamma_{\mathrm{eff}}$.
    \item \textbf{K4:} Matching to $R + R^2/(6M^2)$ after imposing the GR renormalization condition.
\end{itemize}
\end{tcolorbox}

\subsection{K.1 Setup}
We evaluate the one-loop effective action in a Riemann--Cartan background with torsion
decomposed into axial $S_\mu$, trace $T_\mu$, and tensor $q_{\mu\nu\rho}$. The minimal UFE
choice treats torsion at quadratic order, so integrating out torsion is Gaussian and
parameter-free once $\cthree$ fixes the normalization of the coupling.

\subsection{K.2 Operator Form}
After background-field gauge fixing, the quadratic operator for torsion fluctuations can be
written in Laplace type,
\begin{equation}
    \Delta = -\nabla^2 + \mathcal{E}(R,K),
\end{equation}
with $\mathcal{E}$ built from curvature and torsion invariants. The contribution to
$\Gamma_{\mathrm{eff}}$ is $\frac{1}{2}\ln\det\Delta$.

\subsection{K.3 Heat-Kernel Coefficients}
The local part of $\Gamma_{\mathrm{eff}}$ is governed by the heat-kernel expansion
\begin{equation}
    \mathrm{Tr}\,e^{-s\Delta} \sim \sum_{n\ge 0} s^{(n-4)/2} a_n,
\end{equation}
where $a_2$ contains the curvature-squared terms. The torsion-dependent pieces combine into
an effective $R^2$ operator in the low-energy action.

\subsection{K.4 Renormalization Condition}
The GR limit (A8) fixes the renormalization scheme by requiring that the $K\to 0$ limit
reduces to the Einstein--Hilbert action with coefficient $\Mpl^2/2$.

\subsection{K.5 TFPT Closure}
With the TFPT normalization of the torsion sector, the induced $R^2$ coefficient is fixed by
$\cthree$, yielding
\begin{equation}
    \frac{M}{\Mpl}=\sqrt{8\pi}\;\cthree^4,
\end{equation}
which reproduces \eqref{eq:R2scale} without additional parameters.

%=============================================================================
\section{Bounce Perturbations and the Transfer Function}
\label{app:bounce}
%=============================================================================

\begin{tcolorbox}[colback=gray!5, colframe=gray!60, 
    title={\textbf{Status Note}}, fonttitle=\bfseries\small]
This appendix provides a minimal outline; the full numerical transfer-function evaluation
is implemented in the \texttt{bounce\_perturbations} verification module (see \texttt{tfpt-suite/}).
\end{tcolorbox}

\subsection{L.1 Background Solution}
The bounce background is modeled by a smooth scale factor $a(\eta)$ where the torsion-induced
$a^{-6}$ term regularizes $z''/z$ for scalar perturbations. The transition to the $R^2$
inflationary phase is implemented by matching $a(\eta)$ and its derivatives at $\eta=\eta_t$.

\subsection{L.2 Numerical Integration Setup}
Mode functions are initialized in an adiabatic vacuum well before the bounce and integrated
through the bounce and into the $R^2$ era. The numerical scheme records the transfer function
by comparing late-time solutions to the pure $R^2$ baseline.

\subsection{L.3 Transfer Function}
Define the transfer function as
\begin{equation}
    T(k) \equiv \left|\frac{v_k(\eta\gg \eta_t)}{v_k^{(R^2)}(\eta\gg \eta_t)}\right|,
\end{equation}
with $T(k)\to 1$ for modes well inside the horizon during the bounce and deviations confined
to $k\sim k_{\text{bounce}}$.

\TFPTTransferFunctionPlot

